\section{Limitations and Future Work}\label{sec:limits}

\paragraph{Single domain, single language.}
The evaluation domain is one Korean public-sector compensation
regulation. Generalization to other regulatory domains, to English
or other languages, and to private-sector enterprise data is left to
future work. Two of the four LLM families (HCX-32B-Think, Qwen3-235B)
have substantial Korean pre-training; the other two (Llama-4-Scout,
gpt-oss-120b) are predominantly English. Form-sensitivity findings
should be interpreted within this LLM-family mix.

\paragraph{Sample size.}
The 100-question evaluation set is balanced across ten categories
but yields modest per-category power: the
\texttt{out\_of\_scope} category in particular has $n=6$, so the
boundary-inference finding (Section~6.2) is qualitative. Scaling
to 1{,}000+ questions, with explicit balance toward boundary and
multi-key categories, is the highest-priority follow-on.

\paragraph{Production-agent specifics.}
The KG-agent results are tied to one specific production
implementation (LangGraph ReAct with three-attempt retries). A
different agent — for example, a graph-aware retrieval-then-rerank
pipeline like GraphRAG \citep{edge2024graphrag} or a
PageRank-driven retriever like HippoRAG
\citep{gutierrez2024hipporag} — could close some of the
30-percentage-point gap. We do not include such external baselines
in the present paper; doing so is a critical next step before
generalizing the "query-generation bottleneck" claim beyond
ReAct-style agents.

\paragraph{Surface forms covered.}
We test four surface forms (TQL, Cypher, NL, JSON-LD). Other
surface forms — RDF/Turtle, the OpenIE-style triple format used by
some retrieval pipelines, and graph-summary text from GraphRAG —
remain untested. The Pareto-optimality claim for NL is strictly
within the four tested forms.

\paragraph{Schema coverage of in-context dumps.}
The in-context dumps cover the core fact tables (salary, position
pay, differential, cap, wage peak, bonus rate, overseas pay,
starting step, six categories of position) but omit several
production-side entities: career-track entities, allowance entities,
12 supplementary clauses (부칙) with their override relations, and
salary tables for non-comprehensive-planning grades (GA/CL/PO).
The KG agents have full graph access; the in-context tiers do not.
The asymmetry favors the KG agents, which makes the observed in-
context-tier advantage \emph{stronger} than the headline numbers
suggest.

\paragraph{Single non-reasoning configuration.}
All chat LLMs run in non-reasoning mode (no thinking-flag
activation, no chain-of-thought scaffolding). Whether reasoning
modes — particularly for HCX-32B-Think and gpt-oss-120b — would
narrow the KG-agent gap is a natural ablation we leave as future
work.

\paragraph{Closed-book leakage.}
Closed-book accuracy is uniformly $\leq 13\%$, indicating low
pre-training memorization of these specific Korean compensation
facts. However, the four LLMs have heterogeneous training cutoffs
and language coverage, and we do not control for the precise
publication date of training corpora. Replication on data published
\emph{after} all model cutoffs would strengthen the claim that the
RAG accuracy gap is genuinely retrieval, not memorization.

\paragraph{Future work.}
Beyond the items above, we identify three directions: (i) systematic
guide ablation across all four LLMs and three formal surface forms,
(ii) explicit-boundary NL variants with measured lift on
\texttt{out\_of\_scope} and \texttt{negation\_exception}, and (iii)
multi-domain extension covering at least one English-language
regulatory corpus to test cross-language transfer of the
form-sensitivity mechanism.
